\chapter{Dahlberg Step 2: the closing parameter and the converse (Phase D-C)}
\label{chap:dahlberg_step2}

% archon:covers Gluck/Euclidean/DahlbergStep2.lean

This chapter completes Dahlberg's proof of the plane case
(\cref{thm:dahlberg_converse}). Phase D-B (\cref{chap:dahlberg_step1}) produced a
preliminary diffeomorphism $\eta$ with $\kappa\circ\eta=a(1-f)+bf+e$, $0<a<b$,
$\int_0^{2\pi}|e|<C\varepsilon$, and total curvature $I>0$. Step~2 finds a second
adjustment --- a \emph{closing parameter} ranging over the in-tree configuration
disk --- so that the fully reparametrised curvature becomes a non-normalised
curvature function (satisfies (1.2) closure and (1.3) simplicity); the D-A
reduction \cref{thm:realizes_of_non_normalised} then yields the curve.

\medskip\noindent\textbf{Route note (no Möbius disk; Route A node-placing family).}
Dahlberg parametrises the closing adjustment by a Möbius transformation
$g_\beta(w)=(w-\beta)/(1-\bar\beta w)$ over the unit $\beta$-disk and proves the
winding fact via his Propositions 2.2--2.3. We do \emph{not} formalise that family.
Instead we run the existence argument over the in-tree affine configuration disk
\cref{def:configuration_space}, reusing the in-tree winding/degree core
(\cref{lem:error_map_winds_boundary}, \cref{lem:winding_number_c_perturb},
\cref{thm:existence_of_zero}) in place of Dahlberg's Props 2.2--2.3. The bridge
that makes this possible is the keystone
\cref{lem:arclength_error_clean_eq_errorMap}:
$F(z)=c(z)\cdot\mathrm{errorMap}\,a\,b\,\delta\,z$ with a positive scalar
$c(z)=1/\lambda(z)>0$. \emph{Crucially}, the closing family
\cref{def:closing_family} is the \textbf{node-placing arc-length}
reparametrisation, \emph{not} the inclination-tuned \cref{def:align_reparam} of the
positive-case reduction --- the latter does \emph{not} land the arc-length
cumulative tangent angles on $\mathrm{configSpace}$ (this is why the earlier fixed
prefactor keystone $F=\tfrac{a+b}{2}\,\mathrm{errorMap}$ was false). This is the
standing project directive.

\medskip\noindent\textbf{Why node-placing makes $F=c\cdot\mathrm{errorMap}$ true.}
The in-tree error map is the \emph{inclination} closure defect
$\mathrm{errorMap}\,a\,b\,\delta\,z=\int_0^{2\pi}e^{i\theta}/\kappa(\theta)\,d\theta$,
where $\kappa(\theta)$ is the step curvature with jumps at the configuration nodes
$(\theta_1,\theta_2,\theta_3,\theta_4)=\mathrm{configSpace}\,\delta\,(z.\mathrm{re},z.\mathrm{im})$
and values $a,b,a,b,a$ on the five arcs $(0,\theta_1),\dots,(\theta_4,2\pi)$
(\cref{lem:error_vector_of_bicircle}). The arc-length closure defect
$F(z)=\int_0^{2\pi}e^{i\alpha_{K_z}(s)}\,ds$ becomes this same integral under the
substitution $\theta=\alpha_{K_z}(s)$, $d\theta=K_z(s)\,ds$ --- \emph{provided the
cumulative arc-length tangent angle $\alpha_{K_z}$ jumps the step-curvature value
exactly at the nodes $\theta_j$}. Choosing the arc-length interval lengths $L_j$ so
that $\alpha_{K_z}$ advances by exactly $\Delta_j=\theta_j-\theta_{j-1}$ across the
$j$-th clean arc forces this landing, and the Jacobian $1/K_z$ contributes the
constant scalar $c(z)=I_z/2\pi=1/\lambda(z)$.

\medskip
The genuinely new analytic content of this phase is then threefold and is isolated
in \cref{lem:arclength_error_tendsto} (uniform limit
$F^*(\cdot,\varepsilon)\to F$), \cref{lem:total_curvature_ne_zero}
($I^*\neq 0$ for small $\varepsilon$), and \cref{lem:angle_dist_le} (the
$|\alpha-\alpha^*|\le C\varepsilon$ estimate). The simplicity transport
(\cref{lem:simplicity_transport}) and the final assembly
(\cref{thm:dahlberg_converse}) are then short.

\section{The node-placing closing family}

% SOURCE: [Dahlberg], §3, p. 2135, Step 2 (read from references/dahlberg.pdf)
% SOURCE QUOTE: "We now claim that if $\epsilon$ has been chosen small enough,
% then there is a $\beta \in C$, $|\beta| \le \frac{1}{2}$ such that
% $\kappa^* \circ g_\beta$ is a non-normalised curvature function. To see this set
% $k^* = \kappa^* \circ g_\beta$, $k = a(1 - f \circ g_\beta) + b f \circ g_\beta$
% and let $I = \int_0^{2\pi} k\,dt$, $I^* = \int_0^{2\pi} k^*\,dt$,
% $K^* = \frac{2\pi}{I^*} k^*$ and $K = \frac{2\pi}{I} k$."
% --- Route A: a project-bespoke replacement for Dahlberg's Möbius family. It is
% NOT a lemma in Dahlberg; the strategy-auditor confirmed Route A is sound but
% not source-faithful. The fixed-prefactor inclination family (alignReparam) is
% replaced by the node-placing arc-length family below.
\begin{definition}

[Closing reparametrisation family --- node-placing arc-length reparam]
  \label{def:closing_family}
  % Lean (private, not pinned): closingFamily, closingDelta1, closingDelta2, closingDelta3, closingDelta4, closingDelta5, closingLambda, closingLambdaRaw, closingLambdaFloor, closingLen1, closingLen2, closingLen3, closingLen4, closingLen5, closingS1, closingS2, closingS3, closingS4, closingMid1, closingMid2, closingMid3, closingMid4, closingMid5, closingRamp, closingBase, closingHeight, closingDensity
  \uses{def:configuration_space, def:clean_bicircle, def:align_reparam}
  Fix $0<a<b$ and $0<\delta\le\pi/8$. For $z\in\overline{\mathbb{D}}$ write the
  configuration nodes
  \[
    (\theta_1,\theta_2,\theta_3,\theta_4)=\mathrm{configSpace}\,\delta\,(z.\mathrm{re},z.\mathrm{im})
    =\bigl(\tfrac{\pi}{4}+\delta z.\mathrm{re},\,\tfrac{3\pi}{4}+\delta z.\mathrm{im},\,\tfrac{5\pi}{4},\,\tfrac{7\pi}{4}\bigr),
  \]
  with $\theta_0=0$, $\theta_5=2\pi$, and let the five \emph{cumulative-angle}
  increments and clean-curvature values be
  \[
    \Delta_j=\theta_j-\theta_{j-1}\quad(j=1,\dots,5),\qquad
    (\hat\kappa_1,\dots,\hat\kappa_5)=(a,b,a,b,a),
  \]
  matching the five-arc value pattern of \cref{lem:error_vector_of_bicircle}
  (and of $k_{a,b}=a(1-f)+bf$, \cref{def:clean_bicircle}). Put
  \[
    \lambda(z)=\frac1{2\pi}\sum_{j=1}^5\frac{\Delta_j}{\hat\kappa_j}
      =\frac1{2\pi}\Bigl(\frac{\pi+\delta(z.\mathrm{re}-z.\mathrm{im})}{a}
        +\frac{\pi+\delta(z.\mathrm{im}-z.\mathrm{re})}{b}\Bigr)>0,
    \qquad L_j(z)=\frac{\Delta_j}{\lambda(z)\,\hat\kappa_j}>0,
  \]
  so that $\sum_{j=1}^5 L_j(z)=2\pi$ (the two value-$a$ end arcs
  $\Delta_1,\Delta_5$ contribute to the same $1/a$ term). Let the canonical
  clean-bicircle breakpoints be $c_0=0$, $c_1=\tfrac{\pi}{4}$, $c_2=\tfrac{3\pi}{4}$,
  $c_3=\tfrac{5\pi}{4}$, $c_4=\tfrac{7\pi}{4}$, $c_5=2\pi$ (the FIXED jump points of
  $k_{a,b}$, with arc widths $w_j=c_j-c_{j-1}$ equal to
  $\tfrac{\pi}{4},\tfrac{\pi}{2},\tfrac{\pi}{2},\tfrac{\pi}{2},\tfrac{\pi}{4}$ and
  values $a,b,a,b,a$). Define $g_z=\mathrm{closingFamily}\,a\,b\,\delta\,z:\mathbb{R}\to\mathbb{R}$
  to be the $C^1$ orientation-preserving circle reparametrisation
  ($g_z(s+2\pi)=g_z(s)+2\pi$, $g_z'>0$, jointly continuous in $(z,s)$, $g_z(0)=0$)
  that maps the $j$-th \emph{arc-length} interval $[s_{j-1},s_j]$
  ($s_j=\sum_{k\le j}L_k(z)$, so $s_0=0$, $s_5=2\pi$) onto the $j$-th canonical
  clean-bicircle arc $[c_{j-1},c_j]$. Concretely $g_z(s)=\int_0^s w_z(u)\,du$ is the
  running integral of a continuous positive $2\pi$-periodic density built as a
  constant \emph{baseline} plus five \emph{calibrated clamped-tent pulses}, one
  centred on each arc-length interval:
  \[
    w_z=m+\sum_{j=1}^5 h_j\,T_j,\qquad m=\mathrm{closingBase}=\tfrac{\lambda(z)\,a}{2},
    \quad \eta=\mathrm{closingRamp}=\tfrac{\pi a}{20(a+b)},
  \]
  where $T_j=\mathrm{clampTent}\,\eta\,L_j\,C_j$ is the unit-height trapezoidal
  clamped tent centred on the arc midpoint $C_j=s_{j-1}+L_j/2$, with one-period
  support window $[s_{j-1},s_j]$ (the full arc, width $L_j$; value $0$ at the
  endpoints), ramp intervals of width $\eta$ at each end, plateau
  $[s_{j-1}+\eta,\,s_j-\eta]$ of width $L_j-2\eta$, and trapezoidal area
  $\int T_j=L_j-\eta$, and
  $h_j=\mathrm{closingHeight}=\dfrac{w_j-m\,L_j}{\max(\eta,\,L_j-\eta)}
       =\dfrac{w_j-m\,L_j}{L_j-\eta}$
  the calibrated pulse height (the $\max$-clamp is inactive on the disk, where
  $2\eta\le L_j$). The ramp width $\eta$ is $a,b$-dependent and below the
  compact-disk lower bound on the $L_j$ (the fixed $\pi/16$ ramp of
  \cref{def:align_reparam} is NOT reusable, since $L_j(z)$ can be small when $b/a$
  is large). The calibration makes the per-arc integral exactly the canonical
  width, $\int_{s_{j-1}}^{s_j}w_z=m\,L_j+h_j(L_j-\eta)=w_j$ (so the per-arc
  \emph{average} slope is $w_j/L_j(z)$, though the density is not flat on the arc);
  this is the node-mapping \cref{lem:closing_family_nodes}. This is the arc-length
  analogue of \cref{def:align_reparam}; it plays the role of Dahlberg's Möbius
  $\{g_\beta\}$.

  The point of the calibration $L_j=\Delta_j/(\lambda\hat\kappa_j)$ is the
  node-landing \cref{lem:clean_arclength_nodes}: with the clean weight $k_{a,b}$,
  the cumulative normalised arc-length tangent angle satisfies
  $\alpha_{K_z}(s_j)=\theta_j$ --- the configuration nodes, \emph{not} the canonical
  $c_j$.
\end{definition}

\begin{lemma}

\leanok
[The closing family is a running-integral reparametrisation]
  \label{lem:closing_family_eq}
  % Lean (private, not pinned): closingFamily_eq
  \uses{def:closing_family, def:integral_reparam}
  For $0<a<b$, $0<\delta\le\pi/8$ and $z\in\overline{\mathbb{D}}$, the node-placing
  family is the generic running integral of the closing density anchored at $0$:
  \[
    \mathrm{closingFamily}\,a\,b\,\delta\,z
      =\mathrm{integralReparam}\,(\mathrm{closingDensity}\,a\,b\,\delta\,z)\,0 ,
  \]
  i.e.\ $g_z(s)=\int_0^s w_z(u)\,du$ with anchor constant $0$.
\end{lemma}
\begin{proof}
  \uses{def:closing_family, def:integral_reparam}
  \leanok
  Both sides are by definition the map
  $s\mapsto\int_0^s\mathrm{closingDensity}\,a\,b\,\delta\,z$: the generic
  reparametrisation $\mathrm{integralReparam}\,w\,c$ is the running integral of $w$
  from the anchor $c$, and here $c=0$, which is exactly the definition of
  $\mathrm{closingFamily}$.
\end{proof}

% --- New-family API package (the lemmas the prover creates for closingFamily) ---
\begin{lemma}
[The calibration scalar $\lambda(z)$ is positive and continuous]
  \label{lem:closing_lambda_pos}
  % Lean (private, not pinned): closingLambda_pos, continuous_closingLambda, closingLambda_ne, closingLambda_eq_raw, continuous_closingLambdaRaw, closingLambda_le
  \uses{def:closing_family, def:configuration_space}
  For $0<a<b$, $0<\delta\le\pi/8$, the map $z\mapsto\lambda(z)$ is continuous and
  strictly positive on $\overline{\mathbb{D}}$, with uniform bounds
  $0<\lambda_{\min}\le\lambda(z)\le\lambda_{\max}$; consequently the prefactor
  $c(z)=1/\lambda(z)$ is continuous with $0<c_{\min}\le c(z)\le c_{\max}$.
\end{lemma}
\begin{proof}
  \uses{def:configuration_space}
  $\lambda(z)$ is an explicit affine-over-$\{a,b\}$ combination of the configuration
  nodes (\cref{def:configuration_space}), each of which is a continuous (affine)
  function of $z.\mathrm{re},z.\mathrm{im}$; the increments $\Delta_j$ are continuous
  and, for $\delta\le\pi/8$ and $\|z\|\le1$, strictly positive (the gaps of
  $\mathrm{configSpace}$ stay in $[\pi/8,3\pi/8]$ etc.), so $\lambda(z)>0$. On the
  compact disk continuity gives the uniform bounds; positivity of the bounds follows
  since $\lambda$ is continuous and nowhere zero on a compact set. Reciprocation
  gives the same for $c(z)$.
\end{proof}

\begin{lemma}
[Arc-length interval bounds and total]
  \label{lem:closing_len_bounds}
  % Lean (private, not pinned): closingLen_bounds, closingLen_sum, closingLen_bound_aux
  \uses{def:closing_family, lem:closing_lambda_pos}
  For $0<a<b$, $0<\delta\le\pi/8$ and $z\in\overline{\mathbb{D}}$, the five
  arc-length intervals $L_j(z)=\Delta\theta_j/(\lambda(z)\hat\kappa_j)$ satisfy
  $0<L_j$, $2\eta\le L_j$, and $L_j\le 2\pi$ for each $j=1,\dots,5$, and
  $\sum_{j=1}^5 L_j(z)=2\pi$.
\end{lemma}
\begin{proof}
  \uses{def:closing_family, lem:closing_lambda_pos}
  The total identity is the calibration $\lambda=\lambda_{\mathrm{raw}}=
  \tfrac1{2\pi}\sum_j\Delta\theta_j/\hat\kappa_j$, whence
  $\sum_j L_j=\tfrac1\lambda\sum_j\Delta\theta_j/\hat\kappa_j=2\pi$. Each
  $\Delta\theta_j$ is a configuration increment bounded in $[\pi/8,3\pi/8]\cdot$const
  and $\lambda\hat\kappa_j>0$ is uniformly bounded (\cref{lem:closing_lambda_pos}),
  giving the per-interval bounds $0<L_j\le 2\pi$; the clamp-inactivity bound
  $2\eta\le L_j$ holds for the fixed small ramp width $\eta$ since each $L_j$ is
  bounded below by a positive constant on the disk. (A private auxiliary
  $\mathrm{closingLen\_bound\_aux}$ packages the shared per-arc estimate; the
  positivity needs an explicit-argument application to avoid a \texttt{whnf}
  heartbeat blowup.)
\end{proof}

\begin{lemma}
[Basic interface of the node-placing family]
  \label{lem:closing_family_props}
  % Lean (private, not pinned): closingFamily_zero, closingFamily_add_two_pi, closingFamily_two_pi, strictMono_closingFamily, continuous_uncurry_closingFamily, hasDerivAt_closingFamily, closingDensity_pos, closingFamily_slope_bounds, closingRamp_pos, closingDensity_periodic, continuous_closingFamily, continuous_closingBase, continuous_closingDensity_s, continuous_closingLen1, continuous_closingLen2, continuous_closingLen3, continuous_closingLen4, continuous_closingLen5, continuous_closingMid1, continuous_closingMid2, continuous_closingMid3, continuous_closingMid4, continuous_closingMid5, continuous_heightAux, continuous_lenAux, continuous_uncurry_closingDensity, continuous_uncurry_term, closingDensity_period_integral, clampTent_period_integral, comp_closingFamily_periodic
  \uses{def:closing_family, lem:closing_lambda_pos, lem:closing_family_eq}
  The family $g_z=\mathrm{closingFamily}\,a\,b\,\delta\,z$ satisfies: $g_z(0)=0$;
  quasiperiodicity $g_z(s+2\pi)=g_z(s)+2\pi$ (hence $g_z(2\pi)=2\pi$); strict
  monotonicity; everywhere-differentiability $\mathrm{HasDerivAt}\,g_z\,(w_z(s))\,s$
  with $w_z>0$ continuous; joint continuity of $(z,s)\mapsto g_z(s)$; and uniform
  slope/inverse-slope bounds $0<m_0\le w_z(s)\le M_0$ on $\overline{\mathbb{D}}$
  (with $m_0,M_0$ depending only on $a,b,\delta$).
\end{lemma}
\begin{proof}
  \uses{def:closing_family, lem:closing_lambda_pos}
  All clauses follow from the running-integral construction
  $g_z(s)=\int_0^s w_z$, exactly as the corresponding facts for
  \cref{def:align_reparam} (\cref{lem:align_density_props},
  \cref{lem:continuous_uncurry_align_reparam}): $g_z(0)=0$ is the empty integral;
  $w_z$ continuous positive $2\pi$-periodic gives $\mathrm{HasDerivAt}$ (FTC,
  \texttt{integral\_hasDerivAt\_right}), strict monotonicity, and
  $g_z(s+2\pi)=g_z(s)+\int_0^{2\pi}w_z=g_z(s)+2\pi$ (the density integrates to
  $\sum_j w_j=2\pi$). Joint continuity comes from continuity of the plateau slopes
  $w_j/L_j(z)$ in $z$ (\cref{lem:closing_lambda_pos}) and of the ramp profile. The
  slope bounds are the compact-disk extrema of the finitely many plateau slopes and
  ramp slopes.
\end{proof}

\begin{lemma}

[Node-mapping of the closing family onto the canonical clean arcs]
  \label{lem:closing_family_nodes}
  % Lean (private, not pinned): closingFamily_node, closingDensity_arcs, closingDensity_integral_split, closingHeight_mul, half_le_arccos_cos_wide, clampTent_integral_eq_zero_wide
  \uses{def:closing_family, lem:closing_family_props, lem:closing_lambda_pos, lem:closing_len_bounds}
  For $0<a<b$, $0<\delta\le\pi/8$ and $z\in\overline{\mathbb{D}}$, the node-placing
  family sends each arc-length breakpoint to the corresponding canonical
  clean-bicircle breakpoint:
  \[
    \mathrm{closingFamily}\,a\,b\,\delta\,z\,(s_j)=c_j,\qquad j=0,\dots,5,
  \]
  where $s_j=\sum_{k\le j}L_k(z)$ ($s_0=0$, $s_5=2\pi$) and
  $(c_0,\dots,c_5)=(0,\tfrac{\pi}{4},\tfrac{3\pi}{4},\tfrac{5\pi}{4},\tfrac{7\pi}{4},2\pi)$.
  Equivalently the per-arc density integral is
  $\int_{s_{j-1}}^{s_j}w_z=w_j$ (the canonical arc width
  $w_j=c_j-c_{j-1}$). Consequently $g_z$ maps the open arc $(s_{j-1},s_j)$
  strictly monotonically onto $(c_{j-1},c_j)$, so $k_{a,b}\circ g_z=\hat\kappa_j$
  a.e.\ on $(s_{j-1},s_j)$.
\end{lemma}
\begin{proof}
  \uses{def:closing_family, lem:closing_family_props, def:clean_bicircle}
  By \cref{def:closing_family} the density is $w_z=m+\sum_{k=1}^5 h_k\,T_k$, where
  $m=\mathrm{closingBase}=\tfrac{\lambda a}{2}$ is the baseline, $T_k=\mathrm{clampTent}$
  is the unit-height trapezoidal clamped tent $\mathrm{clampTent}\,\eta\,L_k\,C_k$
  centred on the $k$-th arc midpoint $C_k=s_{k-1}+L_k/2$, with one-period support
  window $[s_{k-1},s_k]$ (the full $k$-th arc, value $0$ at the endpoints), ramps
  of width $\eta$, plateau $[s_{k-1}+\eta,s_k-\eta]$ of width $L_k-2\eta$, and
  trapezoidal area $\int T_k=L_k-\eta$; and
  $h_k=\dfrac{w_k-m\,L_k}{\max(\eta,L_k-\eta)}=\dfrac{w_k-m\,L_k}{L_k-\eta}$ the
  calibrated pulse height (the clamp is inactive: $2\eta\le L_k$ on the disk, so
  $\max(\eta,L_k-\eta)=L_k-\eta$). The support windows of distinct $T_k$ meet only
  at shared arc endpoints (where the pulse value is $0$), so on the \emph{open} arc
  $(s_{j-1},s_j)$ only the $j$-th pulse is nonzero and
  \[
    \int_{s_{j-1}}^{s_j}w_z=m\,L_j+h_j\bigl(\textstyle\int_{s_{j-1}}^{s_j}T_j\bigr)
      =m\,L_j+h_j\,(L_j-\eta)=m\,L_j+(w_j-m\,L_j)=w_j .
  \]
  Here the per-arc tent integral $\int_{s_{j-1}}^{s_j}T_j=L_j-\eta$ is the
  trapezoidal \emph{area} (NOT the support width $L_j$): plateau of width
  $L_j-2\eta$ at height $1$ plus two unit-height ramps of base $\eta$ contribute
  $(L_j-2\eta)+2\cdot\tfrac{\eta}{2}=L_j-\eta$. The vanishing of the \emph{other}
  pulses $T_k$ ($k\neq j$) on $(s_{j-1},s_j)$ is the in-tree periodic
  off-support-zero fact; the existing helper assumes support half-width below
  $\pi/2$ (i.e.\ $L<\pi$), but for the node-placing family $L_k$ can exceed $\pi$
  (e.g.\ $L_3\le5\pi/3$ as $b/a\to\infty$), so the periodic copies must be excluded
  for the full range $0<L<2\pi$. This generalised off-support zero is provable from
  \texttt{Real.arccos\_cos} by the case split $y\le\pi$ / $y>\pi$ (a $\sim 30$-line
  helper). Telescoping the running integral,
  $g_z(s_j)=\int_0^{s_j}w_z=\sum_{k\le j}w_k=c_j$. Strict monotonicity
  (\cref{lem:closing_family_props}) upgrades the endpoint equalities to the
  open-arc bijection, and the a.e.\ value $k_{a,b}\circ g_z=\hat\kappa_j$ on
  $(s_{j-1},s_j)$ then follows since $g_z$ maps that open arc into the canonical
  clean-bicircle arc $(c_{j-1},c_j)$, on which $k_{a,b}=a(1-f)+bf$ takes the
  constant value $\hat\kappa_j$ (\cref{def:clean_bicircle}).
\end{proof}

\begin{lemma}

\leanok
[Node-landing: cumulative normalised arc-length angle hits $\mathrm{configSpace}$]
  \label{lem:clean_arclength_nodes}
  % Lean (private, not pinned): cleanArcLength_node_values
  \uses{def:closing_family, def:arclength_norm, lem:clean_prefactor, lem:closing_family_nodes}
  Let $K_z=\mathrm{arcLengthNorm}\,(k_{a,b})\,a\,b\,\delta\,z$ be the normalised clean
  arc-length curvature and $\alpha_{K_z}(s)=\int_0^s K_z$ its cumulative tangent
  angle (\cref{def:dahlberg_angle}). At the arc-length breakpoints
  $s_j=\sum_{k\le j}L_k(z)$,
  \[
    \alpha_{K_z}(s_j)=\theta_j\quad(j=0,\dots,5),\qquad
    (\theta_0,\dots,\theta_5)=(0,\theta_1,\theta_2,\theta_3,\theta_4,2\pi),
  \]
  with $(\theta_1,\dots,\theta_4)=\mathrm{configSpace}\,\delta\,(z.\mathrm{re},z.\mathrm{im})$.
  Moreover $K_z$ is constant $=\lambda(z)\hat\kappa_j$ on each open arc
  $(s_{j-1},s_j)$.
\end{lemma}
\begin{proof}
  \uses{lem:clean_prefactor, lem:closing_family_nodes}
  \leanok
  By \cref{lem:clean_prefactor}, $k_{a,b}$ is a.e.\ the step function with
  value $\hat\kappa_j$ on $(c_{j-1},c_j)$, and $g_z$ maps $(s_{j-1},s_j)$ onto
  $(c_{j-1},c_j)$ (\cref{def:closing_family}), so $k_{a,b}(g_z(s))=\hat\kappa_j$ for
  $s\in(s_{j-1},s_j)$. The clean total curvature is
  $I_z=\int_0^{2\pi}k_{a,b}(g_z)=\sum_j\hat\kappa_j L_j(z)
   =\sum_j\hat\kappa_j\frac{\Delta_j}{\lambda\hat\kappa_j}=\frac1\lambda\sum_j\Delta_j
   =\frac{2\pi}{\lambda}$ (\cref{lem:clean_prefactor}), so the normalising factor is
  $2\pi/I_z=\lambda$ and $K_z=\lambda\,\hat\kappa_j$ on $(s_{j-1},s_j)$. Hence
  \[
    \alpha_{K_z}(s_j)=\int_0^{s_j}K_z=\sum_{k\le j}\lambda\hat\kappa_k L_k(z)
      =\sum_{k\le j}\lambda\hat\kappa_k\frac{\Delta_k}{\lambda\hat\kappa_k}
      =\sum_{k\le j}\Delta_k=\theta_j. \qedhere
  \]
\end{proof}

\begin{definition}

[Normalised arc-length curvature weight]
  \label{def:arclength_norm}
  % Lean (private, not pinned): arcLengthNorm, dahlbergAngle_arcLengthNorm
  \uses{def:closing_family}
  For a curvature weight $g:\mathbb{R}\to\mathbb{R}$, parameters $a,b,\delta$ and
  configuration $z\in\overline{\mathbb{D}}$, the \emph{normalised arc-length weight}
  is the $\tfrac{2\pi}{I_z}$-rescaled reparametrised weight
  \[
    K_z(s)=\frac{2\pi}{I_z}\,(g\circ g_z)(s),\qquad
    I_z=\int_0^{2\pi}(g\circ g_z)(t)\,dt,
  \]
  where $g_z=\mathrm{closingFamily}\,a\,b\,\delta\,z$ (\cref{def:closing_family}).
  This is the common encoding of Dahlberg's normalisations $K_z$ (clean,
  $g=k_{a,b}$) and $K^*_z$ (perturbed, $g=\kappa\circ\eta$); both $F$ and $F^*$ are
  built from it. When $I_z>0$ and $g\ge 0$ a.e., $K_z$ has unit total mass
  $\int_0^{2\pi}K_z=2\pi$.
\end{definition}

\begin{definition}

\leanok
[Perturbed and clean arc-length error maps]
  \label{def:arclength_error_map}
  % Lean (private, not pinned): arcLengthErrorMap
  \uses{def:arclength_norm, def:closing_family, def:dahlberg_angle, def:dahlberg_curve}
  Fix $\kappa,\eta,a,b,e$ from \cref{thm:exists_preliminary_diffeo} at a parameter
  $\varepsilon>0$. For $z\in\overline{\mathbb{D}}$ set
  $k^*_z = (\kappa\circ\eta)\circ g_z$, $k_z = k_{a,b}\circ g_z$, with total
  curvatures $I^*_z,I_z$ and normalisations $K^*_z,K_z$ as above. The
  \emph{perturbed} and \emph{clean} arc-length error maps are the closure defects
  \[
    F^*(z,\varepsilon)=\int_0^{2\pi}e^{i\alpha_{K^*_z}(s)}\,ds=\gamma_{K^*_z}(2\pi),
    \qquad
    F(z)=\int_0^{2\pi}e^{i\alpha_{K_z}(s)}\,ds=\gamma_{K_z}(2\pi),
  \]
  where $\alpha_{K}$ and $\gamma_K$ are the arc-length tangent angle and curve of
  \cref{def:dahlberg_angle,def:dahlberg_curve}. A zero $F^*(z^*,\varepsilon)=0$ is
  exactly condition (1.2) for $K^*_{z^*}$.
\end{definition}

\begin{lemma}

\leanok
[Continuity of the error maps]
  \label{lem:arclength_error_continuous}
  % Lean (private, not pinned): continuous_arcLengthErrorMap
  \uses{def:arclength_error_map, lem:closing_family_props}
  For each fixed $\varepsilon$, $z\mapsto F^*(z,\varepsilon)$ is continuous on
  $\overline{\mathbb{D}}$, and $z\mapsto F(z)$ is continuous.
\end{lemma}
\begin{proof}
  \uses{def:arclength_error_map, lem:closing_family_props, lem:total_curvature_ne_zero}
  \leanok
  The family $g_z$ is jointly continuous in $(z,t)$
  (\cref{lem:closing_family_props}); $k^*_z$ and $I^*_z$ are continuous in $z$
  (composition and parametric integration of continuous functions, as in
  \cref{lem:kappa_error_map_continuous}); $I^*_z>0$ stays away from $0$ by
  \cref{lem:total_curvature_ne_zero}, so $K^*_z$ and then $\alpha_{K^*_z}$,
  $e^{i\alpha_{K^*_z}}$ and the outer integral are continuous in $z$. Likewise for
  $F$, using $I_z=2\pi/\lambda(z)>0$ (\cref{lem:clean_prefactor}).
\end{proof}

\section{The clean error map is a positive multiple of the in-tree \texttt{errorMap}}

The keystone \cref{lem:arclength_error_clean_eq_errorMap} is proved
\emph{arc-by-arc}: the node-landing \cref{lem:clean_arclength_nodes} fixes the
cumulative-angle breakpoints at the configuration nodes $\theta_j$, the prefactor
\cref{lem:clean_prefactor} supplies $c(z)=1/\lambda(z)$, and the general arc-length
arc-integral tool \cref{lem:arclength_arc_integral} evaluates the closure integral
on each of the five constant-curvature arcs. Because the normalised weight $K_z$ is
a step function the global $C^1$ change of variables fails; the arc-by-arc route
avoids it entirely (no $\alpha_{K_z}^{-1}$).

\begin{lemma}
[Clean total curvature and the positive prefactor $c(z)=1/\lambda(z)$]
  \label{lem:clean_prefactor}
  % Lean (private, not pinned): cleanTotalCurvature_eq, cleanPrefactor_pos, cleanBicircle_arcs, clean_arc_ae, clean_arc_int
  \uses{def:closing_family, def:clean_bicircle, lem:closing_lambda_pos, lem:closing_len_bounds, lem:closing_family_nodes}
  For $0<a<b$, $0<\delta\le\pi/8$ and $z\in\overline{\mathbb{D}}$,
  \[
    I_z=\int_0^{2\pi}k_{a,b}(g_z\,t)\,dt=\frac{2\pi}{\lambda(z)},
    \qquad c(z):=\frac{I_z}{2\pi}=\frac1{\lambda(z)}>0 ,
  \]
  with $c$ continuous and bounded between positive constants on
  $\overline{\mathbb{D}}$. (This \emph{replaces} the false fixed value $(a+b)\pi$:
  the clean total curvature is $z$-dependent for the node-placing family.)
\end{lemma}
\begin{proof}
  \uses{lem:closing_lambda_pos, lem:closing_family_nodes}
  The clean weight $k_{a,b}(g_z\,\cdot)$ is a.e.\ the step
  function with value $\hat\kappa_j$ on $(s_{j-1},s_j)$ (since $g_z$ maps that
  interval into the canonical arc $(c_{j-1},c_j)$ of value $\hat\kappa_j$,
  \cref{def:closing_family}). Hence
  $I_z=\sum_{j=1}^5\hat\kappa_j\,(s_j-s_{j-1})=\sum_j\hat\kappa_j L_j(z)
   =\sum_j\hat\kappa_j\frac{\Delta_j}{\lambda(z)\hat\kappa_j}
   =\frac1{\lambda(z)}\sum_j\Delta_j=\frac{2\pi}{\lambda(z)}$. Positivity,
  continuity and the uniform bounds of $c(z)=1/\lambda(z)$ are
  \cref{lem:closing_lambda_pos}.
\end{proof}

\begin{lemma}

\leanok
[Arc-length arc integral on a constant-curvature arc]
  \label{lem:arclength_arc_integral}
  % Lean (private, not pinned): arcLengthArcIntegral
  \uses{def:dahlberg_angle}
  Let $K:\mathbb{R}\to\mathbb{R}$ be \emph{globally interval-integrable}
  ($\forall p\,q,\ \mathrm{IntervalIntegrable}\,K\,\mathrm{volume}\,p\,q$) and
  constant $K\equiv m$ on the open arc $(c,d)$ ($c\le d$) with $m\neq0$, and write
  $\alpha_K(s)=\int_0^s K$ for the arc-length tangent angle (\cref{def:dahlberg_angle}).
  Then
  \[
    \int_c^d e^{\,i\,\alpha_K(s)}\,ds
      =\frac{1}{m}\,(-i)\,\bigl(e^{\,i\,\alpha_K(d)}-e^{\,i\,\alpha_K(c)}\bigr).
  \]
  This is the arc-length analogue of the inclination-variable building block
  $\mathrm{bicircle\_arc\_integral}$ of \cref{lem:error_vector_of_bicircle}. The
  global interval-integrability hypothesis (not merely integrability on $[c,d]$) is
  needed because $\alpha_K$ integrates from $0$, so the split passes through $0$; it
  is met at the use site ($K_z$ is globally interval-integrable).
\end{lemma}
\begin{proof}
  \uses{def:dahlberg_angle}
  \leanok
  For $s\in[c,d]$, $\alpha_K(s)=\alpha_K(c)+\int_c^s K=\alpha_K(c)+m\,(s-c)$ because
  $K=m$ a.e.\ on $(c,d)$ (the two endpoints are a null set). Hence on $[c,d]$ the
  integrand is $e^{i\alpha_K(s)}=e^{i\alpha_K(c)}\,e^{\,i\,m\,(s-c)}$. The map
  $s\mapsto \tfrac{1}{im}\,e^{i\alpha_K(c)}\,e^{\,i\,m\,(s-c)}$ has derivative
  $e^{i\alpha_K(s)}$ on $[c,d]$, so by the fundamental theorem of calculus the
  integral equals
  $\tfrac{1}{im}\,e^{i\alpha_K(c)}\bigl(e^{\,i\,m\,(d-c)}-1\bigr)
   =\tfrac{1}{m}(-i)\bigl(e^{i\alpha_K(d)}-e^{i\alpha_K(c)}\bigr)$
  (using $1/i=-i$ and $\alpha_K(d)=\alpha_K(c)+m(d-c)$). No global $C^1$ structure
  of $\alpha_K$ is used --- only constancy of $K$ on the single arc.
\end{proof}

% NOTE: The earlier statement F = ((a+b)/2)·errorMap was FALSE (the inclination-
% tuned alignReparam does not land arc-length nodes on configSpace). Route A's
% node-placing family (def:closing_family) makes the TRUE identity below hold.
\begin{lemma}

\leanok
[Arc-length closure of the clean bicircle $=$ positive multiple of \texttt{errorMap}]
  \label{lem:arclength_error_clean_eq_errorMap}
  % Lean (private, not pinned): arcLengthError_clean_eq_errorMap
  \uses{def:arclength_error_map, def:error_map, lem:error_vector_of_bicircle, lem:clean_prefactor, lem:arclength_arc_integral, lem:clean_arclength_nodes}
  For $0<a<b$, $0<\delta\le\pi/8$ and $z\in\overline{\mathbb{D}}$, the clean
  arc-length error map $F(z)=\mathrm{arcLengthErrorMap}\,(k_{a,b})\,a\,b\,\delta\,z$
  is a \emph{positive} (configuration-dependent) real multiple of the in-tree error
  map:
  \[
    F(z)=c(z)\cdot\mathrm{errorMap}\,a\,b\,\delta\,z,\qquad c(z)=\frac1{\lambda(z)}>0 .
  \]
  In particular $F(z)=0\iff\mathrm{errorMap}\,a\,b\,\delta\,z=0$, and for any loop in
  $\overline{\mathbb{D}}\setminus\{F=0\}$ the two maps have the same winding number
  around $0$ (positive-scalar-field invariance, \cref{lem:clean_error_winds}).
\end{lemma}
\begin{proof}
  \uses{def:error_map, lem:error_vector_of_bicircle, lem:clean_prefactor, lem:arclength_arc_integral, lem:clean_arclength_nodes}
  \leanok
  Write $K_z=\mathrm{arcLengthNorm}\,(k_{a,b})\,a\,b\,\delta\,z$, so
  $F(z)=\int_0^{2\pi}e^{i\alpha_{K_z}(s)}\,ds$.

  \emph{Arc structure and node landing.} By \cref{lem:clean_arclength_nodes} there
  are breakpoints $0=s_0<s_1<\dots<s_5=2\pi$ with $K_z$ constant $\hat\kappa_j\lambda(z)$
  on each open arc $(s_{j-1},s_j)$ (values $a,b,a,b,a$ scaled by $\lambda$), and the
  cumulative angle landing $\alpha_{K_z}(s_j)=\theta_j$, where $\theta_0=0$,
  $\theta_5=2\pi$ and $(\theta_1,\theta_2,\theta_3,\theta_4)=\mathrm{configSpace}\,\delta\,(z.\mathrm{re},z.\mathrm{im})$.

  \emph{Arc-length split.} Splitting $F(z)$ at $s_1,\dots,s_4$ and applying
  \cref{lem:arclength_arc_integral} with $m=\hat\kappa_j\lambda(z)$ on each arc
  (the global interval-integrability of $K_z$ holds since $K_z$ is a bounded a.e.-step
  function),
  \[
    F(z)=\sum_{j=1}^{5}\frac{1}{\hat\kappa_j\lambda(z)}\,(-i)\,
      \bigl(e^{\,i\,\theta_j}-e^{\,i\,\theta_{j-1}}\bigr)
      =\frac1{\lambda(z)}\sum_{j=1}^{5}\frac{1}{\hat\kappa_j}\,(-i)\,
      \bigl(e^{\,i\,\theta_j}-e^{\,i\,\theta_{j-1}}\bigr).
  \]

  \emph{Match.} By the arc-by-arc evaluation of \cref{lem:error_vector_of_bicircle}
  ($\mathrm{bicircleErrorVector\_eq}$, whose $\mathrm{bicircle\_arc\_integral}$
  building block contributes exactly $\tfrac1{\hat\kappa_j}(-i)(e^{i\theta_j}-e^{i\theta_{j-1}})$
  on the arc $(\theta_{j-1},\theta_j)$ of value $\hat\kappa_j$), the sum
  $\sum_{j}\tfrac{1}{\hat\kappa_j}(-i)(e^{i\theta_j}-e^{i\theta_{j-1}})$ with the same
  nodes $\theta_j$ and the same arc values $a,b,a,b,a$ is precisely
  $\mathrm{bicircleErrorVector}\,a\,b\,\theta_1\theta_2\theta_3\theta_4
   =\mathrm{errorMap}\,a\,b\,\delta\,z$ (\cref{def:error_map}). Therefore
  $F(z)=\tfrac1{\lambda(z)}\,\mathrm{errorMap}\,a\,b\,\delta\,z=c(z)\,\mathrm{errorMap}\,a\,b\,\delta\,z$
  with $c(z)>0$ (\cref{lem:clean_prefactor}). The vanishing equivalence is immediate;
  the equal-winding claim is \cref{lem:clean_error_winds}.
\end{proof}

\section{The three new analytic estimates}

\begin{definition}

\leanok
[Bundled preliminary-diffeomorphism data]
  \label{def:preliminary_diffeo_data}
  % Lean (private, not pinned): PreliminaryDiffeoData
  % DESIGN NOTE (iter-031): error $e$ is interval-integrable, NOT continuous (same
  % reasoning as `thm:exists_preliminary_diffeo`); only the $L^1$ bound is consumed
  % downstream. Lean carries `IntervalIntegrable e volume 0 (2π)`.
  \uses{def:clean_bicircle, thm:exists_preliminary_diffeo}
  At a fixed $\varepsilon>0$, $\mathrm{PreliminaryDiffeoData}\,\kappa\,a\,b\,C\,\varepsilon\,\eta\,e$
  records the output of \cref{thm:exists_preliminary_diffeo}: a $C^1$
  orientation-preserving $\eta$ (through a continuous positive derivative $v$, with
  $\eta(t+2\pi)=\eta(t)+2\pi$); an interval-integrable $2\pi$-periodic error $e$; the
  decomposition $\kappa\circ\eta=a(1-f)+bf+e$ (\cref{def:clean_bicircle}); the $L^1$
  smallness $\int_0^{2\pi}|e|<C\varepsilon$; and the total-curvature estimate
  $\bigl|\int_0^{2\pi}\kappa(\eta(t))\,dt-(a+b)\pi\bigr|<C\varepsilon$. The last
  clause is retained so the consuming estimates obtain $I>0$ and the $I^*$-vs-$I$
  comparison directly.
\end{definition}

% SOURCE: [Dahlberg], §3, p. 2135, "If $\epsilon$ has been chosen small enough, then clearly $I^* \neq 0$" (read from references/dahlberg.pdf)
% SOURCE QUOTE: "If $\epsilon$ has been chosen small enough, then clearly
% $I^* \neq 0$ whenever $|\beta| \le \frac{1}{2}$."
\begin{lemma}

[Total curvature stays nonzero]
  \label{lem:total_curvature_ne_zero}
  % Lean (private, not pinned): totalCurvature_ne_zero, closingFamily_changeOfVar, closingFamily_comp_L1, measurable_cleanBicircle, intervalIntegrable_cleanBicircle_comp
  \uses{def:arclength_error_map, def:preliminary_diffeo_data, lem:closing_family_props}
  There is $\varepsilon_1>0$ such that for all $\varepsilon\in(0,\varepsilon_1)$
  and all $z\in\overline{\mathbb{D}}$, $I^*_z=\int_0^{2\pi}(\kappa\circ\eta)\circ
  g_z>0$ (so in particular $I^*_z\neq 0$ and $K^*_z$ is well defined).
\end{lemma}
\begin{proof}
  \uses{thm:exists_preliminary_diffeo, lem:closing_family_props, lem:clean_prefactor}
  Change variables by the diffeomorphism $g_z$:
  $I^*_z=\int_0^{2\pi}(\kappa\circ\eta)(g_z)\,$; since $g_z$ is an
  orientation-preserving $2\pi$-translation-equivariant diffeomorphism with
  $g_z'=w_z\in[m_0,M_0]$ (\cref{lem:closing_family_props}), $I^*_z=\int_0^{2\pi}(\kappa\circ\eta)(g_z)$.
  Writing $\kappa\circ\eta=k_{a,b}+e$, the clean part contributes
  $I_z=2\pi/\lambda(z)\ge 2\pi/\lambda_{\max}>0$ (\cref{lem:clean_prefactor}) and the
  error contributes at most $\int_0^{2\pi}|e(g_z)|=\int_0^{2\pi}|e|\,(g_z^{-1})'\le
  (\sup(g_z^{-1})')\,C\varepsilon\le \tfrac1{m_0}C\varepsilon$. Choosing
  $\varepsilon_1$ so the error term is below the clean lower bound gives $I^*_z>0$
  uniformly in $z$.
\end{proof}

\begin{lemma}

\leanok
[A $2\pi$-periodic bound on $[0,2\pi]$ holds everywhere]
  \label{lem:abs_le_of_periodic_two_pi}
  % Lean (private, not pinned): abs_le_of_periodic_two_pi
  Let $D:\mathbb{R}\to\mathbb{R}$ be $2\pi$-periodic ($D(s+2\pi)=D(s)$) and suppose
  $|D(s)|\le B$ for all $s\in[0,2\pi]$. Then $|D(s)|\le B$ for every real $s$.
\end{lemma}
\begin{proof}

\leanok
  Reduce $s$ to its representative $s_0\in[0,2\pi)$ modulo $2\pi$, i.e.\
  $s_0=s-k\,2\pi$ for the integer $k$ with $s_0\in[0,2\pi)$. Periodicity gives
  $D(s)=D(s_0)$, and since $s_0\in[0,2\pi]$ the hypothesis bounds $|D(s_0)|\le B$.
  Hence $|D(s)|\le B$.
\end{proof}

% SOURCE: [Dahlberg], §3, p. 2135, "the fact that $\lim_{\epsilon \downarrow 0} F^*(\beta,\epsilon) = F(\beta)$" (read from references/dahlberg.pdf)
% SOURCE QUOTE: "By a standard topology argument and the fact that
% $\lim_{\epsilon \downarrow 0} F^*(\beta,\epsilon) = F(\beta)$ it follows that if
% $\epsilon_0$ has been chosen sufficiently small, there is a $\beta$,
% $|\beta| \le \frac{1}{2}$ such that $F^*(\beta,\epsilon) = 0$."
\begin{lemma}

[Angle estimate $|\alpha-\alpha^*|\le C\varepsilon$]
  \label{lem:angle_dist_le}
  % Lean (private, not pinned): angle_dist_le, dahlbergAngle_arcLengthNorm_add_two_pi, intervalIntegrable_expI_arcLengthNorm, angle_dist_arith
  \uses{def:arclength_error_map, def:preliminary_diffeo_data, lem:total_curvature_ne_zero, lem:clean_prefactor, lem:abs_le_of_periodic_two_pi}
  There is a constant $C'$ (independent of $\varepsilon$ and $z$) such that for
  $\varepsilon$ small, for all $z\in\overline{\mathbb{D}}$ and all $s$,
  $\bigl|\alpha_{K^*_z}(s)-\alpha_{K_z}(s)\bigr|\le C'\varepsilon$.
\end{lemma}
\begin{proof}
  \uses{thm:exists_preliminary_diffeo, lem:total_curvature_ne_zero, lem:clean_prefactor}
  $\alpha_{K^*_z}(s)-\alpha_{K_z}(s)=\int_0^s(K^*_z-K_z)$, with
  $K^*_z=\tfrac{2\pi}{I^*_z}k^*_z$, $K_z=\tfrac{2\pi}{I_z}k_z$ and
  $k^*_z-k_z=e\circ g_z$. Split
  $K^*_z-K_z=\tfrac{2\pi}{I^*_z}(k^*_z-k_z)+\bigl(\tfrac{2\pi}{I^*_z}-\tfrac{2\pi}{I_z}\bigr)k_z$.
  The first term integrates (in absolute value over $[0,2\pi]$) to at most
  $\tfrac{2\pi}{I^*_z}\int|e\circ g_z|=\tfrac{2\pi}{I^*_z}\int|e|(g_z^{-1})'
   \le\tfrac{2\pi}{I^*_z}\tfrac1{m_0}C\varepsilon\le C''\varepsilon$, since $I^*_z$ is
  bounded below (\cref{lem:total_curvature_ne_zero}). For the second, $I_z=2\pi/\lambda(z)$
  is bounded below (\cref{lem:clean_prefactor}), $|I^*_z-I_z|=|\int e\circ g_z|\le
  C''\varepsilon$, so $|\tfrac{2\pi}{I^*_z}-\tfrac{2\pi}{I_z}|\le C'''\varepsilon$ and
  $\int_0^{2\pi}|k_z|=I_z$ is bounded above. Summing over $\int_0^s$ gives
  $\bigl|\alpha_{K^*_z}(s)-\alpha_{K_z}(s)\bigr|\le C'\varepsilon$ for all
  $s\in[0,2\pi]$.

  \emph{Extension to all real $s$ (periodicity reduction).} Both normalised angles
  satisfy $\alpha_{K^*_z}(s+2\pi)=\alpha_{K^*_z}(s)+2\pi$ and
  $\alpha_{K_z}(s+2\pi)=\alpha_{K_z}(s)+2\pi$ (each normalised curvature integrates to
  $2\pi$ over one period, by \cref{lem:total_curvature_ne_zero} and
  \cref{lem:clean_prefactor}). Hence the difference
  $D(s)=\alpha_{K^*_z}(s)-\alpha_{K_z}(s)$ is exactly $2\pi$-periodic, $D(s+2\pi)=D(s)$,
  so the bound $|D|\le C'\varepsilon$ on the fundamental window $[0,2\pi]$ extends to
  every real $s$. (This all-$s$ form is what \cref{lem:simplicity_transport} needs for
  the complement interval $[\tau,t+2\pi]$.)
\end{proof}

\begin{lemma}

[Uniform limit $F^*(\cdot,\varepsilon)\to F$]
  \label{lem:arclength_error_tendsto}
  % Lean (private, not pinned): arcLengthErrorMap_tendsto, norm_expI_sub_le
  \uses{def:arclength_error_map, def:preliminary_diffeo_data, lem:angle_dist_le}
  $\sup_{z\in\overline{\mathbb{D}}}\bigl\|F^*(z,\varepsilon)-F(z)\bigr\|\to 0$ as
  $\varepsilon\downarrow 0$; quantitatively
  $\|F^*(z,\varepsilon)-F(z)\|\le 2\pi\,C'\varepsilon$.
\end{lemma}
\begin{proof}
  \uses{lem:angle_dist_le}
  $F^*(z,\varepsilon)-F(z)=\int_0^{2\pi}(e^{i\alpha_{K^*_z}(s)}-e^{i\alpha_{K_z}(s)})\,ds$.
  Using $|e^{ix}-e^{iy}|\le|x-y|$ and \cref{lem:angle_dist_le},
  $\|F^*(z,\varepsilon)-F(z)\|\le\int_0^{2\pi}|\alpha_{K^*_z}-\alpha_{K_z}|\le 2\pi C'\varepsilon$,
  uniformly in $z$.
\end{proof}

\section{Existence of a closing parameter}

% SOURCE: [Dahlberg], §3, p. 2135, winding argument (read from references/dahlberg.pdf)
% SOURCE QUOTE: "It follows from Propositions 2.1 and 2.2 that $F(\beta) \neq 0$
% whenever $\beta \neq 0, |\beta| < 1$. Noticing that $F(0) = 0$ we see from
% Proposition 2.3 that the restriction of $F$ to a sufficiently small circle
% centred at $0$ has non-vanishing winding number around $0$."
\begin{lemma}

\leanok
[Clean error map: nonzero on the boundary with nonzero winding]
  \label{lem:clean_error_winds}
  % Lean (private, not pinned): cleanError_winds_boundary
  \uses{lem:arclength_error_clean_eq_errorMap, lem:closing_lambda_pos, lem:error_map_winds_boundary, lem:winding_number_c_pos_scalar_field}
  For $0<a<b$ and $0<\delta\le\pi/8$, the clean error map $F$ is nonzero on the unit
  circle $\partial\mathbb{D}$ and its winding number around $0$ along
  $\partial\mathbb{D}$ is nonzero.
\end{lemma}
\begin{proof}
  \uses{lem:arclength_error_clean_eq_errorMap, lem:closing_lambda_pos, lem:error_map_winds_boundary, lem:winding_number_c_pos_scalar_field}
  \leanok
  By \cref{lem:arclength_error_clean_eq_errorMap}, $F(z)=c(z)\cdot\mathrm{errorMap}\,a\,b\,\delta\,z$
  with $c(z)=1/\lambda(z)>0$ continuous on $\partial\mathbb{D}$
  (\cref{lem:closing_lambda_pos}). The in-tree \cref{lem:error_map_winds_boundary}
  ($\mathrm{errorMap\_winding\_eq\_one}$) states $\mathrm{errorMap}$ is nonvanishing
  on $\partial\mathbb{D}$ with boundary winding $-1$. Multiplication by a
  \emph{positive scalar field} $c(z)>0$ changes neither the nonvanishing nor the
  winding number: the loops $t\mapsto\mathrm{errorMap}(t)$ and
  $t\mapsto c(t)\,\mathrm{errorMap}(t)$ are homotopic through
  $(u,t)\mapsto\bigl((1-u)+u\,c(t)\bigr)\,\mathrm{errorMap}(t)$, whose scalar factor
  $(1-u)+u\,c(t)>0$ is positive throughout (so the homotopy avoids $0$), and a
  positive-scalar reparametrisation of a loop preserves its winding number. Hence $F$
  is nonvanishing on $\partial\mathbb{D}$ with winding $-1\neq0$. (This is the
  in-tree breakpoint-disk version of Dahlberg's Prop 2.3, with the unit circle in
  place of his "sufficiently small circle".)
\end{proof}

% SOURCE: [Dahlberg], §3, p. 2135, "By a standard topology argument" (read from references/dahlberg.pdf)
% SOURCE QUOTE: "By a standard topology argument and the fact that
% $\lim_{\epsilon \downarrow 0} F^*(\beta,\epsilon) = F(\beta)$ it follows that if
% $\epsilon_0$ has been chosen sufficiently small, there is a $\beta$,
% $|\beta| \le \frac{1}{2}$ such that $F^*(\beta,\epsilon) = 0$."
\begin{theorem}

\leanok
[Existence of a closing parameter]
  \label{thm:exists_closing_param}
  % Lean: private `exists_closingParam` in Gluck/Euclidean/DahlbergStep2.lean (private declarations are not pinned)
  \uses{def:preliminary_diffeo_data, lem:clean_error_winds, lem:arclength_error_tendsto, lem:arclength_error_continuous, lem:winding_number_c_perturb, thm:existence_of_zero}
  Assume $\kappa$ is continuous (the hypothesis $\mathrm{Continuous}\,\kappa$ is
  carried in the Lean signature, supplied at the use site by
  \cref{def:mixed_sign_four_vertex}). There is $\varepsilon_0>0$ such that for every
  $\varepsilon\in(0,\varepsilon_0)$ there is $z^*\in\mathbb{D}$ (interior) with
  $F^*(z^*,\varepsilon)=0$. Hence
  $K^*_{z^*}=\tfrac{2\pi}{I^*_{z^*}}(\kappa\circ\eta)\circ g_{z^*}$ satisfies the
  closure condition (1.2).
\end{theorem}
\begin{proof}
  \uses{lem:clean_error_winds, lem:arclength_error_tendsto, lem:winding_number_c_perturb, thm:existence_of_zero}
  \leanok
  Since $\kappa$ is continuous and $\eta$ is $C^1$ (hence continuous, from
  \cref{def:preliminary_diffeo_data}), the perturbed weight $\kappa\circ\eta$ is
  continuous, so \cref{lem:arclength_error_continuous} applies to both $F$ and
  $F^*(\cdot,\varepsilon)$. By \cref{lem:clean_error_winds}, $F$ is nonvanishing on
  $\partial\mathbb{D}$ with a positive margin $\mu=\min_{|z|=1}\|F(z)\|>0$
  (continuity on the compact circle, \cref{lem:arclength_error_continuous}) and
  nonzero boundary winding. By
  \cref{lem:arclength_error_tendsto} choose $\varepsilon_0$ so that for
  $\varepsilon<\varepsilon_0$, $\sup_z\|F^*(z,\varepsilon)-F(z)\|<\mu$. Then
  $F^*(\cdot,\varepsilon)$ is also nonvanishing on $\partial\mathbb{D}$, and by the
  perturbation invariance of the winding number
  (\cref{lem:winding_number_c_perturb}, applied as in the positive-case
  \cref{lem:reduction_justified}) its boundary winding equals that of $F$, hence is
  nonzero. The planar degree principle \cref{thm:existence_of_zero}
  ($\mathrm{exists\_zero\_of\_boundary\_winding}$) then produces an interior zero
  $z^*\in\mathbb{D}$ with $F^*(z^*,\varepsilon)=0$. Finally
  $F^*(z^*,\varepsilon)=\gamma_{K^*_{z^*}}(2\pi)=0$ is condition (1.2).
\end{proof}

\section{Simplicity transport and the converse}

\begin{lemma}

\leanok
[Chord of the reconstructed curve $=$ integral of its unit tangent]
  \label{lem:dahlberg_curve_sub}
  % Lean: private `dahlbergCurve_sub` in Gluck/Euclidean/DahlbergStep2.lean (private declarations are not pinned)
  \uses{def:dahlberg_curve, def:dahlberg_angle}
  Let $K:\mathbb{R}\to\mathbb{R}$ have unit tangent $s\mapsto e^{i\alpha_K(s)}$
  interval-integrable on every interval (where $\alpha_K$ is the arc-length tangent
  angle of \cref{def:dahlberg_angle}). Then for all $u,w$,
  \[
    \mathrm{dahlbergCurve}\,K\,(w)-\mathrm{dahlbergCurve}\,K\,(u)
      =\int_u^w e^{\,i\,\alpha_K(s)}\,ds .
  \]
\end{lemma}
\begin{proof}
  \uses{def:dahlberg_curve, def:dahlberg_angle}
  \leanok
  By \cref{def:dahlberg_curve}, $\mathrm{dahlbergCurve}\,K\,(t)=\int_0^t
  e^{i\alpha_K(s)}\,ds$. Subtracting the values at $w$ and $u$ and using additivity
  of the integral over the adjacent intervals $[0,u]$ and $[u,w]$
  ($\int_0^w=\int_0^u+\int_u^w$, valid by the interval-integrability hypothesis)
  gives $\int_0^w-\int_0^u=\int_u^w e^{i\alpha_K(s)}\,ds$.
\end{proof}

\begin{lemma}

\leanok
[Two-sided bounds on the normalised clean-bicircle curvature]
  \label{lem:clean_bicircle_arclength_norm_bounds}
  % Lean: private `cleanBicircle_arcLengthNorm_bounds` in Gluck/Euclidean/DahlbergStep2.lean (private declarations are not pinned)
  \uses{def:arclength_norm, def:clean_bicircle, lem:closing_lambda_pos, lem:closing_len_bounds}
  For $0<a<b$, $0<\delta\le\pi/8$ and $\|z\|\le1$, the normalised clean-bicircle
  curvature is pinned between two positive constants: for all $s$,
  \[
    \tfrac14\bigl(\tfrac1a+\tfrac1b\bigr)\,a
      \;\le\;\mathrm{arcLengthNorm}\,(k_{a,b})\,a\,b\,\delta\,z\,(s)
      \;\le\;\tfrac58\bigl(\tfrac1a+\tfrac1b\bigr)\,b .
  \]
\end{lemma}
\begin{proof}
  \uses{def:arclength_norm, def:clean_bicircle, lem:closing_lambda_pos, lem:closing_len_bounds}
  \leanok
  By \cref{def:arclength_norm} and the clean total curvature $I_z=2\pi/\lambda(z)$,
  the normalised weight equals $\lambda(z)\,(k_{a,b}\circ g_z)(s)$, where
  $\lambda(z)=\mathrm{closingLambda}\,a\,b\,\delta\,z>0$
  (\cref{lem:closing_lambda_pos}). The bicircle $k_{a,b}$ takes values in $[a,b]$
  (\cref{def:clean_bicircle}), and the $\mathrm{closingLambda}$ bounds
  (\cref{lem:closing_lambda_pos}, \cref{lem:closing_len_bounds}) pin $\lambda(z)$ into
  $[\tfrac14(\tfrac1a+\tfrac1b),\,\tfrac58(\tfrac1a+\tfrac1b)]$. Multiplying the two
  one-sided estimates gives the stated two-sided bound, uniformly in $z$ and $s$.
\end{proof}

% SOURCE: [Dahlberg], §3, p. 2135, simplicity transport (read from references/dahlberg.pdf)
% SOURCE QUOTE: "Since the curve $\gamma(t) = \int_0^t e^{i\alpha(u)} du$ is simple
% and closed and $|\alpha - \alpha^*| \le C\epsilon$, it follows that
% $\int_t^\tau e^{i\alpha^*(s)} ds \neq 0$ whenever $0 < t < \tau < 2\pi$, which
% completes the proof of the theorem."
\begin{lemma}

\leanok
[Clean chord integrals are bounded away from $0$ on inclination-span-$\le\pi$ arcs]
  \label{lem:clean_chord_margin}
  % Lean: private `clean_chord_margin` in Gluck/Euclidean/DahlbergStep2.lean (private declarations are not pinned)
  % DESIGN NOTE (iter-035): the margin is keyed to the INCLINATION span
  % Λ = α_{K_z}(τ) − α_{K_z}(t), NOT to arc-length (an arc-length-short arc can
  % still sweep large inclination). It does NOT assume clean closure F(z)=0 (the
  % clean curve is NOT closed at the perturbed zero z*). The margin is UNIFORM in
  % z (z* depends on ε, so ε must be chosen below the margin before z* is fixed).
  % DESIGN NOTE (iter-038): the endpoint range is 0 ≤ t < τ ≤ 4π (NOT 2π), so that
  % lem:simplicity_transport can apply the margin on the COMPLEMENT interval
  % [τ, t+2π] ⊆ [0,4π] in its Λ>π case (where 0 ≤ t < τ < 2π ⇒ t+2π < 4π). The
  % compactness proof is unchanged — the set is just the larger compact box [0,4π].
  \uses{def:closing_family, lem:closing_lambda_pos, lem:closing_len_bounds, lem:clean_prefactor, lem:clean_bicircle_arclength_norm_bounds, lem:dahlberg_curve_sub}
  There is a margin $m>0$ (depending only on $a,b,\delta$) such that for every
  $z\in\overline{\mathbb{D}}$, writing $\alpha=\alpha_{K_z}$ for the clean normalised
  curvature $K_z=\mathrm{arcLengthNorm}\,(k_{a,b})\,a\,b\,\delta\,z$ (positive, since
  $a,b>0$),
  \[
    m\,(\tau-t)\;\le\;\bigl\|\textstyle\int_t^\tau e^{i\alpha(s)}\,ds\bigr\|
    \qquad\text{for all } 0\le t<\tau\le 4\pi \text{ with } \alpha(\tau)-\alpha(t)\le\pi .
  \]
\end{lemma}
\begin{proof}
  \uses{def:closing_family, lem:closing_lambda_pos, lem:closing_len_bounds, lem:clean_prefactor}
  \leanok
  The landed proof is elementary and uses no compactness; it is keyed to a uniform
  slope bound plus the cosine positivity of the \emph{middle angular half} of the arc.

  \emph{Uniform slope bounds.} Write $K_z=\lambda(z)\,k_{a,b}(g_z\,\cdot)$ with
  $\lambda(z)=2\pi/I_z$. Since $\lambda(z)$ is bounded above and below
  (\cref{lem:closing_lambda_pos}, \cref{lem:closing_len_bounds}) and
  $k_{a,b}\in[a,b]$ a.e.\ (\cref{lem:clean_prefactor}), there are constants
  $0<m_K\le M_K$, depending only on $a,b,\delta$, with $m_K\le K_z\le M_K$ for every
  $z\in\overline{\mathbb D}$. Hence $\alpha(s)=\int_0^s K_z$ is continuous, strictly
  increasing, with $m_K(y-x)\le\alpha(y)-\alpha(x)\le M_K(y-x)$.

  \emph{Middle-half projection.} Project the chord onto the angular-midpoint direction
  $e^{i\psi}$, $\psi=\tfrac12(\alpha(t)+\alpha(\tau))$:
  $\|\int_t^\tau e^{i\alpha}\|\ge\mathrm{Re}\bigl(e^{-i\psi}\int_t^\tau e^{i\alpha}\bigr)
  =\int_t^\tau\cos(\alpha(s)-\psi)\,ds$. Writing $\Lambda=\alpha(\tau)-\alpha(t)\le\pi$,
  on all of $[t,\tau]$ we have $|\alpha(s)-\psi|\le\Lambda/2\le\pi/2$, so the integrand
  is $\ge0$. By IVT pick $s_1,s_2\in[t,\tau]$ with $\alpha(s_1)=\alpha(t)+\Lambda/4$ and
  $\alpha(s_2)=\alpha(t)+3\Lambda/4$; on $[s_1,s_2]$, $|\alpha(s)-\psi|\le\Lambda/4\le\pi/4$,
  so $\cos(\alpha(s)-\psi)\ge\cos(\pi/4)$.

  \emph{Length of the middle half.} From $\Lambda/2=\alpha(s_2)-\alpha(s_1)\le M_K(s_2-s_1)$
  and $\Lambda\ge m_K(\tau-t)$ we get $s_2-s_1\ge\Lambda/(2M_K)\ge m_K(\tau-t)/(2M_K)$.
  Dropping the nonnegative tails outside $[s_1,s_2]$,
  \[
    \Bigl\|\textstyle\int_t^\tau e^{i\alpha}\Bigr\|\ge\int_{s_1}^{s_2}\cos(\alpha-\psi)
    \ge\cos(\tfrac\pi4)\,(s_2-s_1)\ge m\,(\tau-t),\qquad m=\frac{m_K\cos(\pi/4)}{2M_K}>0,
  \]
  uniform in $z$. This is exactly the half-circle limit ($\Lambda=\pi$, chord/arc $=2/\pi$)
  that a naive whole-arc midpoint bound $(\tau-t)\cos(\Lambda/2)$ misses.
\end{proof}

\begin{lemma}

\leanok
[Simplicity of the perturbed curve (condition (1.3))]
  \label{lem:simplicity_transport}
  % Lean: private `simplicity_transport` in Gluck/Euclidean/DahlbergStep2.lean (private declarations are not pinned)
  \uses{def:preliminary_diffeo_data, lem:clean_chord_margin, lem:angle_dist_le, thm:exists_closing_param, lem:dahlberg_curve_sub}
  For $\varepsilon$ small enough, the closing parameter $z^*$ of
  \cref{thm:exists_closing_param} ($F^*(z^*,\varepsilon)=0$) gives a $K^*_{z^*}$
  satisfying condition (1.3):
  $\int_t^\tau e^{i\alpha_{K^*_{z^*}}(s)}\,ds\neq 0$ for all $0\le t<\tau<2\pi$.
\end{lemma}
\begin{proof}
  \uses{lem:clean_chord_margin, lem:angle_dist_le, thm:exists_closing_param}
  \leanok
  Write $\alpha=\alpha_{K_{z^*}}$ (clean, at the same $z^*$) and $\alpha^*=\alpha_{K^*_{z^*}}$
  (perturbed). Fix $\varepsilon$ with $C'\varepsilon<m$, where $m$ is the uniform clean
  margin of \cref{lem:clean_chord_margin} and $C'$ the angle constant of
  \cref{lem:angle_dist_le}; note the threshold is independent of $z^*$ (both $m$ and $C'$
  are uniform in $z$), so it may be chosen before $z^*$ is produced. The transport
  estimate over any arc $[u,v]$ is, by \cref{lem:angle_dist_le} and $|e^{ix}-e^{iy}|\le|x-y|$,
  \[
    \bigl\|\textstyle\int_u^v e^{i\alpha^*}-\int_u^v e^{i\alpha}\bigr\|
    \le\int_u^v|\alpha^*-\alpha|\le(v-u)\,C'\varepsilon . \tag{$\dagger$}
  \]
  Let $\Lambda=\alpha(\tau)-\alpha(t)$ be the clean inclination span of $[t,\tau]$;
  since $\alpha(2\pi)-\alpha(0)=2\pi$ exactly (normalisation), the complementary span is
  $2\pi-\Lambda$, and exactly one of $\Lambda,2\pi-\Lambda$ is $\le\pi$.

  \emph{Case $\Lambda\le\pi$.} By \cref{lem:clean_chord_margin},
  $\|\int_t^\tau e^{i\alpha}\|\ge m(\tau-t)$, so by $(\dagger)$,
  $\|\int_t^\tau e^{i\alpha^*}\|\ge(m-C'\varepsilon)(\tau-t)>0$.

  \emph{Case $\Lambda>\pi$.} The complement arc $[\tau,t+2\pi]$ has clean span
  $\alpha(t+2\pi)-\alpha(\tau)=2\pi-\Lambda<\pi$, so by \cref{lem:clean_chord_margin}
  $\|\int_\tau^{t+2\pi}e^{i\alpha}\|\ge m(t+2\pi-\tau)>0$, and by $(\dagger)$
  $\|\int_\tau^{t+2\pi}e^{i\alpha^*}\|\ge(m-C'\varepsilon)(t+2\pi-\tau)>0$. Finally, the
  \emph{perturbed} curve is closed ($F^*(z^*,\varepsilon)=\int_0^{2\pi}e^{i\alpha^*}=0$),
  and $e^{i\alpha^*}$ is $2\pi$-periodic ($\alpha^*(s+2\pi)=\alpha^*(s)+2\pi$), so
  $\int_t^\tau e^{i\alpha^*}=-\int_\tau^{t+2\pi}e^{i\alpha^*}$; the two have equal norm,
  hence $\|\int_t^\tau e^{i\alpha^*}\|>0$. In both cases the perturbed chord is nonzero,
  which is (1.3).
\end{proof}

% SOURCE: [Dahlberg], §1, p. 2131, Theorem 1.1 (read from references/dahlberg.pdf)
% SOURCE QUOTE: "Theorem 1.1. Suppose $\kappa : T \to R$ is not a constant and
% assume that $\kappa$ has at least four critical points ordered counterclockwise
% $p_i \in T$, $i = 1,\dots,4$, such that $p_1,p_3$ are local maxima and
% $p_2,p_4$ are local minima with $\kappa(p_1) > 0, \kappa(p_3) > 0$ and
% $\max(\kappa(p_2),\kappa(p_4)) < \min(\kappa(p_1),\kappa(p_3))$. If in addition
% $\kappa$ is smooth, then $\kappa$ is the curvature function of a smooth, simple
% and closed curve."
\begin{theorem}

\leanok
[Dahlberg converse to the four vertex theorem (Theorem 1.1)]
  \label{thm:dahlberg_converse}
  \lean{Gluck.dahlberg_converse}
  \uses{def:mixed_sign_four_vertex, thm:exists_preliminary_diffeo, def:preliminary_diffeo_data, thm:exists_closing_param, lem:total_curvature_ne_zero, lem:simplicity_transport, lem:exists_c1_circle_inverse, lem:closing_family_props, thm:realizes_of_non_normalised, def:non_normalised_curvature}
  \textit{Source: Dahlberg, §1, Theorem 1.1.}
  Let $\kappa:\mathbb{R}\to\mathbb{R}$ be continuous, $2\pi$-periodic and
  non-constant, satisfying the mixed-sign four-vertex condition
  (\cref{def:mixed_sign_four_vertex}). Then there is a simple closed curve realizing
  $\kappa$: $\exists\,\gamma,\ \mathrm{IsSimpleClosed}\,\gamma\,\wedge\,
       \mathrm{RealizesCurvature}\,\gamma\,\kappa$.
\end{theorem}
\begin{proof}
  \uses{thm:exists_preliminary_diffeo, thm:exists_closing_param, lem:simplicity_transport, lem:total_curvature_ne_zero, lem:exists_c1_circle_inverse, lem:closing_family_props, thm:realizes_of_non_normalised}
  \leanok
  From \cref{def:mixed_sign_four_vertex}, $\kappa$ is continuous and $2\pi$-periodic.
  Fix $\delta=\pi/8$. Apply \cref{thm:exists_preliminary_diffeo} to obtain $0<a<b$
  and, for a parameter $\varepsilon$ to be fixed below, the preliminary
  diffeomorphism $\eta$ (with continuous positive derivative $v_\eta$ and
  $\eta(t+2\pi)=\eta(t)+2\pi$, \cref{def:preliminary_diffeo_data}) with
  $\kappa\circ\eta=a(1-f)+bf+e$, $\int_0^{2\pi}|e|<C\varepsilon$. Shrink $\varepsilon$
  below $\min(\varepsilon_0,\varepsilon_1)$ and below the simplicity threshold of
  \cref{lem:simplicity_transport}. By \cref{thm:exists_closing_param} (whose
  $\mathrm{Continuous}\,\kappa$ hypothesis is supplied by
  \cref{def:mixed_sign_four_vertex}) there is an interior $z^*$ with
  $F^*(z^*,\varepsilon)=0$, i.e.\ $K^*_{z^*}$ satisfies (1.2); feeding that same zero
  to \cref{lem:simplicity_transport} gives (1.3); and \cref{lem:total_curvature_ne_zero}
  gives the total curvature $I^*_{z^*}=\int_0^{2\pi}(\kappa\circ\eta)\circ g_{z^*}>0$.

  \emph{The composite reparametrisation and its inverse.} Set
  $\varphi=\eta\circ g_{z^*}$. By the chain rule
  $\mathrm{HasDerivAt}\,\varphi\,\bigl(v_\eta(g_{z^*}(s))\cdot w_{z^*}(s)\bigr)\,s$,
  where $w_{z^*}=g_{z^*}'$ is continuous and strictly positive
  (\cref{lem:closing_family_props}); the product $v_\varphi=
  (v_\eta\circ g_{z^*})\cdot w_{z^*}$ is therefore continuous and strictly positive,
  so $\varphi$ is $C^1$ with $\varphi'>0$. Quasiperiodicity composes:
  $\varphi(s+2\pi)=\eta(g_{z^*}(s)+2\pi)=\eta(g_{z^*}(s))+2\pi=\varphi(s)+2\pi$. By
  \cref{lem:exists_c1_circle_inverse} applied to $\varphi$ (with derivative
  $v_\varphi$) there is a $C^1$ orientation-preserving quasiperiodic two-sided
  inverse $\psi$ with $\psi'=1/v_\varphi\circ\psi>0$.

  \emph{Assembly.} We have $\kappa\circ\varphi=(\kappa\circ\eta)\circ g_{z^*}=k^*_{z^*}$,
  and $K^*_{z^*}=\tfrac{2\pi}{I^*_{z^*}}\kappa\circ\varphi$ with $I^*_{z^*}>0$ satisfies
  (1.2) and (1.3); thus $\kappa\circ\varphi$ is a non-normalised curvature function
  (\cref{def:non_normalised_curvature}) with positive total curvature. Feeding
  $\kappa$, $\varphi$, $\psi$ and all the $C^1$/positivity/quasiperiodicity/inverse
  data to \cref{thm:realizes_of_non_normalised} yields a simple closed curve $\gamma$
  with $\mathrm{RealizesCurvature}\,\gamma\,\kappa$, completing the proof. (This is the
  mixed-sign analogue of \cref{thm:gluck_converse}, stated through the same
  realization predicate per the unification directive.)
\end{proof}

\begin{theorem}

\leanok
[Gluck converse, non-constant strictly-positive case]
  \label{thm:gluck_converse_nonconstant}
  \lean{Gluck.gluck_converse_nonconstant}
  \uses{thm:dahlberg_converse, lem:mixed_sign_four_vertex_of_is_curvature_function}
  Let $\kappa:\mathbb{R}\to\mathbb{R}$ be a strictly positive curvature function
  ($\mathrm{IsCurvatureFunction}\,\kappa$) satisfying the non-constant four-vertex
  condition. Then $\kappa$ is realized by a simple closed curve:
  $\exists\,\gamma,\ \mathrm{IsSimpleClosed}\,\gamma\,\wedge\,
       \mathrm{RealizesCurvature}\,\gamma\,\kappa$.
\end{theorem}
\begin{proof}
  \uses{thm:dahlberg_converse, lem:mixed_sign_four_vertex_of_is_curvature_function}
  \leanok
  A strictly positive four-vertex $\kappa$ satisfies the mixed-sign four-vertex
  hypothesis of \cref{def:mixed_sign_four_vertex}, by
  \cref{lem:mixed_sign_four_vertex_of_is_curvature_function}. Applying
  \cref{thm:dahlberg_converse} to this hypothesis produces the simple closed curve
  realizing $\kappa$. (The omitted constant case is a round circle, handled
  separately by \cref{thm:gluck_converse}.)
\end{proof}

\noindent\textbf{Phase status.} \cref{thm:dahlberg_converse} is the project's
extended goal: the converse to the four vertex theorem \emph{without} global
positivity (positivity only at the two maxima), stated through the same
$\mathrm{IsSimpleClosed}\wedge\mathrm{RealizesCurvature}$ conclusion as the positive
case. It consumes Phase D-A (\cref{chap:arclength}), Phase D-B
(\cref{chap:dahlberg_step1}), and the in-tree winding/degree and chord-integral
machinery, with the genuinely new analytic work confined to the node-placing family
package (\cref{def:closing_family}, \cref{lem:clean_arclength_nodes}) and the
estimates \cref{lem:total_curvature_ne_zero,lem:angle_dist_le,lem:arclength_error_tendsto}.
